\documentclass[10pt,a4paper]{amsart}
\usepackage[margin=19mm]{geometry}
\usepackage{amsmath,amssymb,mathtools}
\usepackage[scr=rsfs,cal=euler]{mathalfa}
\usepackage{xfrac}
\usepackage[unicode,hidelinks,bookmarks=false]{hyperref}
\newcommand{\ie}{i.e.\ }
\newcommand{\cf}{cf.\ }
\newcommand{\resp}{respectively\ }
\newcommand{\Subsection}{\subsection}
\providecommand{\nobreakdash}{\nobreak\hbox{-}}
\setlength{\emergencystretch}{2em}
\multlinegap0pt
\usepackage{bm}
\usepackage{bbm}
\usepackage{dsfont}
\usepackage{xspace}
\usepackage{cancel}
\usepackage{mathtools}
\usepackage{amsthm}
\makeatletter
\@ifundefined{theorem}{\newtheorem{theorem}{Theorem}}{}
\@ifundefined{lemma}{\newtheorem{lemma}[theorem]{Lemma}}{}
\makeatother
\title[Failure of Weak Approximation in Adjoint Groups]{Failure of Weak Approximation in Adjoint Groups}
\author{Chayansudha Biswas}
\date{Source: arXiv:2604.14420v1 (CC BY 4.0)}
\begin{document}
\maketitle

\noindent\emph{Source and licence.} Failure of Weak Approximation in Adjoint Groups. Version arXiv:2604.14420v1 (CC BY 4.0), distributed under the Creative Commons Attribution 4.0 International licence, as recorded in the arXiv licence metadata for this exact version and shown on the abstract page https://arxiv.org/abs/2604.14420v1. Full rights notice: licenses/arxiv-2604.14420-CC-BY-4.0.txt.\par\smallskip

\noindent\textit{Editorial scope.} Counted unit: the Lemma whose source label is \texttt{not-pure-quat-norm} in the article (source0.tex lines 126--128, source label \texttt{not-pure-quat-norm}), delivered together with the article's own complete original proof of it (source0.tex lines 129--149, one \texttt{proof} environment of 21 lines). No mathematics is written, repaired or re-derived: statement and proof are the article's own text line for line. Required context retained uncounted: ab-vals (lines 110--112). Every definition, display and label of those lines is retained unchanged. Omitted from this package and not redistributed: 1--109, 113--125, 150--240. Nothing inside a retained excerpt is omitted or altered: every retained body is delivered exactly as the article's lines are written, so no omission receipt of any kind accompanies this package. Citations: the retained excerpts carry no citation command at all, so no bibliography entry of the article is transcribed and no reference is redistributed. Reference adaptation: none was needed and none was made; every reference inside the delivered unit resolves inside it, so no unresolved reference or citation is produced. Printed slips kept literal: no printed word, symbol, hypothesis or number of the counted statement or of its proof was altered. Layout and class: the article's own class is \texttt{amsart} with packages including amsmath, amsfonts, tikz-cd, amsthm, amssymb, hyperref, mathtools, mathrsfs, wasysym, cite; the wrapper is the lane's pinned \texttt{amsart} editorial wrapper at 10pt on A4, which restates the theorem-like environments the retained text needs and adds the article's own macro definitions that the retained text needs (none), with extra wrapper packages (bm, bbm, dsfont, xspace, cancel, mathtools). The delivered numbering is the unit's own. Assets: no retained span contains a figure environment, a graphics-inclusion command, a PDF-page inclusion, a table or a TikZ picture; no image file is copied into this package, the package declares no asset dependency and no image file is redistributed with it. Licence: version arXiv:2604.14420v1 of this article is distributed under the Creative Commons Attribution 4.0 International licence, as recorded in the arXiv licence metadata for this exact version and shown on the abstract page https://arxiv.org/abs/2604.14420v1. A full rights notice is delivered at licenses/arxiv-2604.14420-CC-BY-4.0.txt. Characters: every retained excerpt is pure ASCII, so the delivered engine, XeLaTeX with its default Latin Modern text font, reports no missing character.
\begin{lemma}\label{ab-vals}
    Given $v\in \Sigma$, there exist $a,b\in F$ with $Q\simeq (a,b)_F$ such that $v(b)=1$ and $v(a)=0$. 
\end{lemma}

\begin{lemma}\label{not-pure-quat-norm}
    The element $d\in F^\times$ is not the norm of a pure quaternion in $Q$.
\end{lemma}

\begin{proof}
    Fix a valuation $v\in \Sigma$ with uniformizer $\pi$ such that the residue field $\kappa_v$ has characteristic different from $2$. By Lemma \ref{ab-vals}, there exist $a,b\in F$ with $Q\simeq (a,b)_F$ and such that $v(a) = 0$ and $v(b)=1$.
    Suppose $d = \operatorname{Nrd}(xI + yJ + zK) = -ax^2-by^2+abz^2$ for some $x,y,z\in F$. By multiplying throughout by a suitable $n$-th power of $\pi^2$, we obtain an equation of the form
    \begin{equation}\label{normform}
        D = -aX^2-bY^2 + abZ^2,
    \end{equation} where $D = d\pi^{2n}$, and $X = x\pi^{n},Y = y\pi^{n}$ and $Z = z\pi^{n}$ are in $\mathcal{O}_v$ and at least one of them is not in $\mathfrak{m}_v$. 
    
    Reducing the equation modulo $\mathfrak{m}_v$, we obtain $D\equiv -aX^2(\operatorname{mod}\mathfrak{m}_v)$. First assume $X\in\mathcal{O}_v^\times$. Since $d$ is a square in $F_v$, so is $D$, and let $C\in F_v$ be such that $C^2 = D$. Then we find $(iC/X)^2\equiv a(\operatorname{mod}\mathfrak{m}_v)$, which would mean that $a$ is a square modulo $\mathfrak{m}_v$, which contradicts the fact that $Q_{F_v}$ is not split. (Here, as before, $i$ is a square root of $-1$ in $F$.)

    So $v(X)>0$, and either $Y$ or $Z$ has zero valuation. Suppose it is $Y$. The case $Z\in\mathcal{O}_v^\times$ is similar. Since $v(b) = 1$, let $b = b_0\pi$ for $b_0\in\mathcal{O}_v^\times$. The congruence $D\equiv -aX^2\operatorname{mod}\mathfrak{m}_v$ gives $D\equiv X^2\equiv 0\operatorname{mod}\mathfrak{m}_v$. Let $D = D_0\pi^k, X= X_0\pi^j$ for $D_0,X_0\in\mathcal{O}_v^\times$ and $k,j\geq 1$. On the other hand, since $D = d\pi^{2n}$, one has $k = v(D) = v(d) + 2n$, which is even because $d$ is a square in $F_v$ (and hence $v(d)\in 2\mathbb{Z}$). Therefore, $k\geq 2$. Substituting these in (\ref{normform}), we obtain:\begin{equation}
        D_0\pi^k = -aX_0^2\pi^{2j} - b_0\pi Y^2 + ab_0\pi Z^2.
    \end{equation}
    Dividing throughout by $\pi$, we get
    \begin{equation}
        D_0\pi^{k-1} = -aX_0^2\pi^{2j-1} - b_0Y^2 + ab_0Z^2;
    \end{equation}
    and reducing this modulo $\mathfrak{m}_v$ gives:\begin{equation}
        b_0Y^2\equiv ab_0Z^2\operatorname{mod}\mathfrak{m}_v.
    \end{equation}
    Since $b_0\in\mathcal{O}_v^\times$, we can multiply both sides by $b_0^{-1}$ to obtain $Y^2\equiv aZ^2\operatorname{mod}\mathfrak{m}_v$. Since $v(Y) = 0$ by assumption and since $v(a)=0$, $Z\not\equiv 0\operatorname{mod}\mathfrak{m}_v$ i.e. $Z\in\mathcal{O}_v^\times$ and so we find $(Y/Z)^2\equiv a\operatorname{mod}\mathfrak{m}_v$, which would mean that $a$ is a square modulo $\mathfrak{m}_v$, again a contradiction.
\end{proof}

\end{document}
